% Provenance (never printed): NSI terminale "Bases de donnees: modele relationnel (relation, attribut, domaine,
% cle primaire, cle etrangere, schema relationnel); contraintes d'integrite (domaine, relation, reference);
% redondance et anomalies; services d'un SGBD; langage SQL (SELECT, FROM, WHERE, JOIN, DISTINCT, ORDER BY, fonctions
% d'agregation sans GROUP BY; INSERT, UPDATE, DELETE)". See sources/book2/BRIEFS.md.
\chapter{Relational Databases and SQL}\label{ch:g12:relational-databases-sql}

The town library once kept its loans in one wide spreadsheet: a row per loan, with the title and the author of the
book, the name of the member, the street and the postcode. It was easy to read. Then a member moved, and her new
address had to be typed into eight \omterm{def:g12:relational-databases-sql:relation}{rows}; one was missed, and the library had two addresses for her. A book whose
loans were all deleted at the end of the year disappeared from the catalogue with them.

This chapter splits such a table into several, linked by keys, so that each fact is written once. The tables and
their rules form a relational \omterm{def:g12:relational-databases-sql:database}{database}, and a language, \omterm{def:g12:relational-databases-sql:sql}{SQL}, asks it questions and changes it. A \omterm{def:g12:relational-databases-sql:database}{database} system
then keeps the tables safe on disk, lets several programs use them at once, and refuses changes that break the rules.

\begin{recall}
A data table is a list of records with the same fields, its columns named by descriptors; the \emph{domain} of a
column is the set of values it may take (chapter \ref{ch:g10:data-tables}). A join combines two tables on a shared
field, and a \omterm{def:g10:data-tables:missing}{missing value} is an empty cell (the same chapter). The town of Brenmoor has its people, streets, books,
members and loans in tables (Book 1, grade 8).
\end{recall}

\section{Why one wide table fails}\label{sec:g12:relational-databases-sql:flat}

The program rebuilds the wide table from the town's tables, and looks at what changing it would cost.

\omcode{code/grade-12/09-relational-databases-sql/listings/flat_table.py}{1}{30}{The library as one wide table: one
row per loan.}\label{lst:g12:relational-databases-sql:flat}

\omoutput{code/grade-12/09-relational-databases-sql/outputs/flat_table.txt}{One fact written many times, and three
kinds of trouble.}\label{out:g12:relational-databases-sql:flat}

\begin{definition}[Redundancy, anomalies]\label{def:g12:relational-databases-sql:anomaly}
A table has \emph{redundancy}\index{redundancy} when the same fact is stored in several places. It then suffers
\emph{anomalies}\index{anomaly (insertion, deletion, update)}: an \emph{update anomaly}, when changing one fact
requires changing several \omterm{def:g12:relational-databases-sql:relation}{rows}, and missing one leaves the table contradicting itself; a \emph{deletion anomaly},
when deleting some \omterm{def:g12:relational-databases-sql:relation}{rows} loses a fact that had nothing to do with them; an \emph{insertion anomaly}, when a fact cannot
be stored until some unrelated fact exists.
\end{definition}

\begin{omfigure}
\begin{tikzpicture}[font=\scriptsize, x=1mm, y=-5mm, cell/.style={draw=black!30, minimum height=5mm, inner sep=1.5pt,
    anchor=west}]
  \foreach \h/\x/\w in {loan/0/10, title/10/30, author/40/24, member/64/26, street/90/22, postcode/112/16}
    \node[cell, minimum width=\w mm, fill=black!8, font=\scriptsize\bfseries] at (\x,0) {\h};
  \foreach \l/\t/\a/\m/\s/\p/\fb/\fm [count=\r] in {
      2/Seven Forest/Elif Ashworth/Bea Stoneford/Bridge Road/B-100/omThm!12/omDef!14,
      53/The Silent Forest/Nils Mossham/Bea Stoneford/Bridge Road/B-100/white/omDef!14,
      63/Seven Forest/Elif Ashworth/Hana Brookwick/Market Road/B-110/omThm!12/white,
      117/The Silent Machine/Zoe Stoneton/Bea Stoneford/Bridge Road/B-100/white/omDef!14} {
    \node[cell, minimum width=10mm] at (0,\r) {\l};
    \node[cell, minimum width=30mm, fill=\fb] at (10,\r) {\t};
    \node[cell, minimum width=24mm, fill=\fb] at (40,\r) {\a};
    \node[cell, minimum width=26mm, fill=\fm] at (64,\r) {\m};
    \node[cell, minimum width=22mm, fill=\fm] at (90,\r) {\s};
    \node[cell, minimum width=16mm, fill=\fm] at (112,\r) {\p};
  }
\end{tikzpicture}

{\small Four \omterm{def:g12:relational-databases-sql:relation}{rows} of the wide table built by \code{flat_table.py}. The shaded cells repeat a fact written in
another row: the author of \emph{Seven Forest} twice here (nine times in the whole table), the address of Bea
Stoneford three times here (eight in all). Each repetition is a place where an update can be forgotten.}
\end{omfigure}

The cure is to write each fact once, in a table about that kind of thing, and to let the loans point to books and
members instead of copying them.

\section{The relational model}\label{sec:g12:relational-databases-sql:model}

\begin{definition}[Database, database management system]\label{def:g12:relational-databases-sql:database}
A \emph{database}\index{database} is a collection of data organised to be stored, queried and changed by programs. A
\emph{database management system}\index{database management system} is the program that stores a database, answers
the queries of other programs, and applies their changes.
\end{definition}

\begin{definition}[Relation, attribute, row]\label{def:g12:relational-databases-sql:relation}
In the relational model, a \emph{relation}\index{relation} is a table: a set of \emph{rows}\index{row (relation)}, all
with the same \emph{attributes}\index{attribute (relation)}, the named columns. Each attribute has a domain. A relation
is a set: its rows have no order, and no two rows are equal.
\end{definition}

\begin{definition}[Primary key, foreign key]\label{def:g12:relational-databases-sql:keys}
A \emph{primary key}\index{primary key} of a \omterm{def:g12:relational-databases-sql:relation}{relation} is an attribute, or a group of \omterm{def:g12:relational-databases-sql:relation}{attributes}, whose values are
different in every row and chosen to identify the \omterm{def:g12:relational-databases-sql:relation}{rows}. A \emph{foreign key}\index{foreign key} is an attribute whose
values must be values of the primary key of another \omterm{def:g12:relational-databases-sql:relation}{relation}: it points to a row of that \omterm{def:g12:relational-databases-sql:relation}{relation}.
\end{definition}

\begin{definition}[Relational schema, integrity constraint]\label{def:g12:relational-databases-sql:schema}
The \emph{relational schema}\index{relational schema} of a \omterm{def:g12:relational-databases-sql:database}{database} lists its \omterm{def:g12:relational-databases-sql:relation}{relations}, with for each one its
\omterm{def:g12:relational-databases-sql:relation}{attributes}, their domains, its \omterm{def:g12:relational-databases-sql:keys}{primary key} and its \omterm{def:g12:relational-databases-sql:keys}{foreign keys}. An \emph{integrity constraint}\index{integrity
constraint} is a rule every state of the \omterm{def:g12:relational-databases-sql:database}{database} must satisfy. Three kinds come with the schema: \emph{domain}
constraints (each value lies in its attribute's domain), \emph{\omterm{def:g12:relational-databases-sql:relation}{relation}} constraints (no two \omterm{def:g12:relational-databases-sql:relation}{rows} share a \omterm{def:g12:relational-databases-sql:keys}{primary
key}) and \emph{reference} constraints (each \omterm{def:g12:relational-databases-sql:keys}{foreign key} points to a row that exists).
\end{definition}

The schema is the \emph{structure} of the \omterm{def:g12:relational-databases-sql:database}{database}; the \omterm{def:g12:relational-databases-sql:relation}{rows} present at a given moment are its \emph{content}. The
structure changes rarely, by design; the content changes with every loan.

\begin{omfigure}
\begin{tikzpicture}[font=\scriptsize, x=1cm, y=-0.42cm,
    att/.style={draw=black!40, minimum width=22mm, minimum height=0.42cm, inner sep=1pt, anchor=center,
      font=\scriptsize\ttfamily},
    head/.style={att, fill=omDef!18, font=\scriptsize\bfseries}]
  \foreach \t/\x/\list in {streets/0/{street_id,name,postcode},
      people/3.1/{person_id,first_name,last_name,birth_year,street_id},
      members/6.2/{member_id,person_id,joined},
      loans/9.3/{loan_id,book_id,member_id,out_date,due_date,return_date},
      books/12.4/{book_id,title,author,year,genre,pages}} {
    \node[head] at (\x,0) {\t};
    \foreach \a [count=\k] in \list {
      \ifnum\k=1 \node[att] (\t-\k) at (\x,\k) {\underline{\detokenize\expandafter{\a}}};
      \else \node[att] (\t-\k) at (\x,\k) {\detokenize\expandafter{\a}}; \fi
    }
  }
  \draw[->, thick, omThm] (people-5.west) -- ++(-0.2,0) |- (streets-1.east);
  \draw[->, thick, omThm] (members-2.west) -- ++(-0.2,0) |- (people-1.east);
  \draw[->, thick, omThm] (loans-3.west) -- ++(-0.2,0) |- (members-1.east);
  \draw[->, thick, omThm] (loans-2.east) -- ++(0.2,0) |- (books-1.west);
\end{tikzpicture}

{\small Part of the town's schema. \omterm{def:g12:relational-databases-sql:keys}{Primary keys} are underlined; each arrow goes from a \omterm{def:g12:relational-databases-sql:keys}{foreign key} to the \omterm{def:g12:relational-databases-sql:relation}{relation}
whose \omterm{def:g12:relational-databases-sql:keys}{primary key} it names: a loan points to a book and to a member, a member to a person, a person to a street.
Each fact now has one place: an address in \code{streets}, an author in \code{books}.}
\end{omfigure}

In \omterm{def:g12:relational-databases-sql:sql}{SQL}, the schema is written as \code{CREATE TABLE} statements. The small library below declares its constraints;
the runner turns on the checking of \omterm{def:g12:relational-databases-sql:keys}{foreign keys}, which SQLite leaves off unless asked.

\omcode{code/grade-12/09-relational-databases-sql/listings/schema.sql}{1}{24}{A small library: keys, \omterm{def:g12:relational-databases-sql:keys}{foreign keys},
and two domain checks.}\label{lst:g12:relational-databases-sql:schema}

\omoutput{code/grade-12/09-relational-databases-sql/outputs/create.txt}{Four \omterm{def:g12:relational-databases-sql:relation}{rows} accepted, then one change refused
for each kind of constraint: domain, \omterm{def:g12:relational-databases-sql:relation}{relation}, reference twice.}\label{out:g12:relational-databases-sql:create}

\section{Querying}\label{sec:g12:relational-databases-sql:query}

\begin{definition}[Query, SQL]\label{def:g12:relational-databases-sql:sql}
A \emph{query}\index{query} is a request for data, answered with a \omterm{def:g12:relational-databases-sql:relation}{relation}. \emph{SQL}\index{SQL} (Structured Query
Language) is the language in which relational \omterm{def:g12:relational-databases-sql:database}{databases} are queried and changed. A query names the columns wanted,
the \omterm{def:g12:relational-databases-sql:relation}{relations} they come from and a condition on the \omterm{def:g12:relational-databases-sql:relation}{rows}:
\code{SELECT columns FROM relation WHERE condition}.
\end{definition}

The files below are run by the book's runner, \code{sqlrun.py}, which loads the town's tables (\code{-- load: town}),
prints each statement after \code{sql>}, and then its answer.

\omoutput{code/grade-12/09-relational-databases-sql/outputs/select_where.txt}{Selecting \omterm{def:g12:relational-databases-sql:relation}{rows} with \code{WHERE}:
comparisons, \code{AND}, and \code{LIKE} with \texttt{\%} for any
letters.}\label{out:g12:relational-databases-sql:select}

\code{DISTINCT} removes repeated \omterm{def:g12:relational-databases-sql:relation}{rows} from an answer; \code{ORDER BY} sorts it, by several columns if needed,
\code{DESC} for decreasing order; \code{LIMIT} keeps the first \omterm{def:g12:relational-databases-sql:relation}{rows}.

\omoutput{code/grade-12/09-relational-databases-sql/outputs/distinct_order.txt}{The eight genres, one author's books
from the newest, and the three longest books.}\label{out:g12:relational-databases-sql:order}

An aggregate function turns a whole selection into one value: \code{COUNT}, \code{SUM}, \code{AVG}, \code{MIN},
\code{MAX}. \code{COUNT(*)} counts \omterm{def:g12:relational-databases-sql:relation}{rows}; \code{COUNT(column)} counts the values that are not missing.

\omoutput{code/grade-12/09-relational-databases-sql/outputs/aggregates.txt}{One row per \omterm{def:g12:relational-databases-sql:sql}{query}: $400$ loans, of which
$378$ have a return date.}\label{out:g12:relational-databases-sql:aggregates}

A \omterm{def:g10:data-tables:missing}{missing value} is written \code{NULL} in \omterm{def:g12:relational-databases-sql:sql}{SQL}. It is not a value like the others: any comparison with it, even
\code{NULL = NULL}, is neither true nor false, so a \code{WHERE} keeps no row on it. The test is \code{IS NULL}.

\omoutput{code/grade-12/09-relational-databases-sql/outputs/still_out.txt}{$22$ loans still out; the comparison
\code{= NULL} finds none.}\label{out:g12:relational-databases-sql:null}

The titles of the books are in \code{books}, the dates in \code{loans}. A join puts together each loan and the book its
\code{book_id} points to. Long names are shortened by \code{AS}.

\omoutput{code/grade-12/09-relational-databases-sql/outputs/joins.txt}{A join of two \omterm{def:g12:relational-databases-sql:relation}{relations}, then of four: the
oldest loans still out, with the names of the borrowers.}\label{out:g12:relational-databases-sql:joins}

\begin{omfigure}
\begin{tikzpicture}[font=\scriptsize, x=1mm, y=-5mm, cell/.style={draw=black!40, minimum height=5mm, inner sep=1.5pt,
    anchor=west}]
  \node[font=\footnotesize\bfseries] at (14,-1.2) {loans};
  \node[font=\footnotesize\bfseries] at (79,-1.2) {books};
  \foreach \h/\x/\w in {loan\_id/0/12, book\_id/12/14} \node[cell, minimum width=\w mm, fill=black!8] at (\x,0) {\h};
  \foreach \l/\b [count=\r] in {1/59, 2/73, 3/28} {
    \node[cell, minimum width=12mm] at (0,\r) {\l};
    \node[cell, minimum width=14mm, fill=omThm!12] (lb\r) at (12,\r) {\b};
  }
  \foreach \h/\x/\w in {book\_id/60/14, title/74/30} \node[cell, minimum width=\w mm, fill=black!8] at (\x,0) {\h};
  \foreach \b/\t [count=\r] in {28/Under the Bridge, 59/The Hidden Valley, 73/Seven Forest} {
    \node[cell, minimum width=14mm, fill=omThm!12] (bb\b) at (60,\r) {\b};
    \node[cell, minimum width=30mm] at (74,\r) {\t};
  }
  \draw[omThm, thick, ->] (lb1.east) -- (bb59.west);
  \draw[omThm, thick, ->] (lb2.east) -- (bb73.west);
  \draw[omThm, thick, ->] (lb3.east) -- (bb28.west);
\end{tikzpicture}

{\small \code{loans JOIN books ON loans.book_id = books.book_id}: each loan is matched with the row of \code{books}
whose key equals its \code{book_id}, and the two \omterm{def:g12:relational-databases-sql:relation}{rows} are written side by side as one row of the answer. Without the
\code{ON} condition, every loan would be paired with every book.}
\end{omfigure}

\section{Changing the data}\label{sec:g12:relational-databases-sql:changes}

\code{INSERT INTO} adds \omterm{def:g12:relational-databases-sql:relation}{rows}, \code{UPDATE ... SET ... WHERE} changes the \omterm{def:g12:relational-databases-sql:relation}{rows} that satisfy a condition, and
\code{DELETE FROM ... WHERE} removes them. The \omterm{def:g12:relational-databases-sql:database}{database} system checks every change against the constraints and refuses
the ones that would break them.

\omoutput{code/grade-12/09-relational-databases-sql/outputs/changes.txt}{A book added and removed, a due date moved,
eight books changed at once; deleting a member who has loans, and a loan of a book that does not exist, are
refused.}\label{out:g12:relational-databases-sql:changes}

An \code{UPDATE} or a \code{DELETE} without \code{WHERE} applies to every row of the \omterm{def:g12:relational-databases-sql:relation}{relation}. That is rarely what was
meant, and a \omterm{def:g12:relational-databases-sql:database}{database} system does not ask for confirmation.

\section{What a database system guarantees}\label{sec:g12:relational-databases-sql:dbms}

A \omterm{def:g12:relational-databases-sql:database}{database management system} offers more than a place to keep tables. It gives \emph{persistence}: a change it has
accepted is on disk and survives a crash or a power cut. It allows \emph{concurrent access}: many programs, on many
machines, read and change the same \omterm{def:g12:relational-databases-sql:database}{database} at once, and each change is applied entirely or not at all, never half.
It answers queries \emph{efficiently}: it keeps indexes, sorted structures over chosen columns, and finds \omterm{def:g12:relational-databases-sql:relation}{rows} without
reading every one. It \emph{controls access}: each user may read or change only what the administrator allowed.

\begin{omfigure}
\begin{tikzpicture}[font=\footnotesize, box/.style={draw, rounded corners=2pt, minimum width=30mm, minimum height=7mm,
    align=center}]
  \node[box, fill=omMeth!12] (p) at (0,0) {a program\\ \scriptsize Python, a web server};
  \node[box, fill=omDef!12] (s) at (0,-1.6) {SQL statements\\ \scriptsize and their answers};
  \node[box, fill=omThm!12, minimum width=52mm] (d) at (0,-3.2) {database management system\\ \scriptsize checks,
    indexes, several users, access rights};
  \node[box, fill=black!6] (f) at (0,-4.8) {files on disk};
  \draw[<->, thick] (p) -- (s); \draw[<->, thick] (s) -- (d); \draw[<->, thick] (d) -- (f);
  \node[box, fill=omMeth!12] (q) at (4.2,0) {another program};
  \draw[<->, thick] (q) |- (s.east);
\end{tikzpicture}

{\small The layers. Programs never touch the files of the \omterm{def:g12:relational-databases-sql:database}{database}: they send \omterm{def:g12:relational-databases-sql:sql}{SQL} and receive answers. The system in
between applies the constraints and lets several programs work on the same data.}
\end{omfigure}

From Python, the module \code{sqlite3} opens a \omterm{def:g12:relational-databases-sql:database}{database}, here in memory, and sends it \omterm{def:g12:relational-databases-sql:sql}{SQL}. The values of a \omterm{def:g12:relational-databases-sql:sql}{query}
are not pasted into its text: they are passed separately, and each \code{?} of the \omterm{def:g12:relational-databases-sql:sql}{query} takes one.

\omcode{code/grade-12/09-relational-databases-sql/listings/from_python.py}{1}{29}{The town's tables from Python, and
a \omterm{def:g12:relational-databases-sql:sql}{query} with two parameters.}\label{lst:g12:relational-databases-sql:python}

\omoutput{code/grade-12/09-relational-databases-sql/outputs/from_python.txt}{Two loans since September by members
named Stoneford.}\label{out:g12:relational-databases-sql:python}

\begin{definition}[SQL injection]\label{def:g12:relational-databases-sql:injection}
An \emph{SQL injection}\index{SQL injection} is an attack on a program that builds the text of a \omterm{def:g12:relational-databases-sql:sql}{query} by pasting
text supplied by a user: the text closes the quotation and adds \omterm{def:g12:relational-databases-sql:sql}{SQL} of the attacker's choice, which the \omterm{def:g12:relational-databases-sql:database}{database} runs.
\end{definition}

\omcode{code/grade-12/09-relational-databases-sql/listings/injection.py}{1}{25}{A \omterm{def:g12:relational-databases-sql:sql}{query} built by pasting, attacked on a
toy table; then the same \omterm{def:g12:relational-databases-sql:sql}{query} with a parameter.}\label{lst:g12:relational-databases-sql:injection}

\omoutput{code/grade-12/09-relational-databases-sql/outputs/injection.txt}{The attack turns the condition into one
that is always true and reads every note; with a parameter, it is only a strange
name.}\label{out:g12:relational-databases-sql:injection}

\begin{safety}
Never build the text of an \omterm{def:g12:relational-databases-sql:sql}{SQL} \omterm{def:g12:relational-databases-sql:sql}{query} by pasting text that came from outside the program: a form, a file, a web
address. Pass every value as a parameter. The attack above was run on a table the program had just created in memory;
trying it on someone else's system is an offence.
\end{safety}

\begin{history}[Codd, 1970]
Edgar F.\ Codd proposed the relational model in 1970, in an article of the \emph{Communications of the ACM} titled
``A relational model of data for large shared data banks''. Programs, he argued, should ask for data by its values,
not by where it is stored.
% ledger: g12:codd-1970
\end{history}

\section{Exercises}\label{sec:g12:relational-databases-sql:exercises}

\begin{exercise}[$\star$]\label{exo:g12:relational-databases-sql:1}
The \omterm{def:g12:relational-databases-sql:relation}{relation} \code{pets(name, kind, age)} holds the \omterm{def:g12:relational-databases-sql:relation}{rows} (Rex, dog, 4), (Tom, cat, 2), (Kiki, bird, 1),
(Max, dog, 9). What does \texttt{SELECT name FROM pets WHERE kind = \textquotesingle{}dog\textquotesingle{} AND age > 5} return?
\end{exercise}

\begin{exercise}[$\star$]\label{exo:g12:relational-databases-sql:2}
A \omterm{def:g12:relational-databases-sql:relation}{relation} \code{seats(train, coach, seat, passenger)} records the seats booked on a day's trains. Which \omterm{def:g12:relational-databases-sql:relation}{attributes}
form a \omterm{def:g12:relational-databases-sql:keys}{primary key}?
\end{exercise}

\begin{exercise}[$\star$]\label{exo:g12:relational-databases-sql:3}
Write a \omterm{def:g12:relational-databases-sql:sql}{query} that counts the different authors of the town library. (The answer is $23$.)
\end{exercise}

\begin{exercise}[$\star$]\label{exo:g12:relational-databases-sql:4}
Write a \omterm{def:g12:relational-databases-sql:sql}{query} that lists the titles and years of the books by Caio Lakeby, oldest first.
\end{exercise}

\begin{exercise}[$\star$]\label{exo:g12:relational-databases-sql:5}
A student writes \code{SELECT COUNT(*) FROM loans, books}, forgetting the join condition. What does it count?
\end{exercise}

\begin{exercise}[$\star\star$]\label{exo:g12:relational-databases-sql:6}
Write a \omterm{def:g12:relational-databases-sql:sql}{query} that lists, once each, the first and last names of the members who have a book still out.
\end{exercise}

\begin{exercise}[$\star\star$]\label{exo:g12:relational-databases-sql:7}
How many \omterm{def:g12:relational-databases-sql:relation}{rows} has \code{loans JOIN books ON loans.book_id = books.book_id}? Why exactly that many?
\end{exercise}

\begin{exercise}[$\star\star$]\label{exo:g12:relational-databases-sql:8}
Why does \code{WHERE return_date = NULL} keep no row, even for the loans still out? Write the right condition.
\end{exercise}

\begin{exercise}[$\star\star$]\label{exo:g12:relational-databases-sql:9}
Write an \code{UPDATE} that gives every loan still out $14$ more days. SQLite's \texttt{date(d, \textquotesingle{}+14 days\textquotesingle{})} adds days to
a date. How many \omterm{def:g12:relational-databases-sql:relation}{rows} change?
\end{exercise}

\begin{exercise}[$\star\star$]\label{exo:g12:relational-databases-sql:10}
With \omterm{def:g12:relational-databases-sql:keys}{foreign keys} checked, which of these are refused on the town's tables, and why: deleting street $1$, loan $1$,
person $1$, member $1$?
\end{exercise}

\begin{exercise}[$\star\star$]\label{exo:g12:relational-databases-sql:11}
A school keeps \code{marks(student, class, teacher, teacher_phone, subject, mark)}, a row per mark. Give an example
of each of the three \omterm{def:g12:relational-databases-sql:anomaly}{anomalies}, and a schema without them.
\end{exercise}

\begin{exercise}[$\star\star\star$]\label{exo:g12:relational-databases-sql:12}
Write the schema of a school timetable: teachers, rooms, and lessons (day, hour, room, teacher, subject), with keys
and constraints so that a room holds one lesson at a time, a teacher gives one lesson at a time, and lessons are on
weekdays between $8$ and $17$.
\end{exercise}

\begin{exercise}[$\star\star\star$]\label{exo:g12:relational-databases-sql:13}
Write \code{books_between(conn, first, last)}, which returns the titles of the books published between two years,
with a parameterised \omterm{def:g12:relational-databases-sql:sql}{query}, and a test for it.
\end{exercise}

\begin{exercise}[$\star\star\star$]\label{exo:g12:relational-databases-sql:14}
Prove that every \omterm{def:g12:relational-databases-sql:relation}{relation} has at least one key: a group of \omterm{def:g12:relational-databases-sql:relation}{attributes} whose values differ in every row.
\end{exercise}

\begin{exercise}[$\star\star\star$]\label{exo:g12:relational-databases-sql:15}
A program runs \texttt{conn.execute(f"DELETE FROM notes WHERE owner = \textquotesingle{}\{name\}\textquotesingle{}")} with a name typed by a user. What name
deletes every note? Rewrite the line safely.
\end{exercise}

\section{Problem: The Overdue Books of Brenmoor}\label{sec:g12:relational-databases-sql:problem}

\begin{problem}[{Weekend problem --- which books were late on the first of September?}]\label{pb:g12:relational-databases-sql:1}
The library wants a picture of its late books on 1 September 2025: which loans were overdue that day, who held them,
where those members live, and how many days of lateness had piled up.

\omcode{code/grade-12/09-relational-databases-sql/listings/overdue.sql}{1}{22}{Four queries about the loans overdue on
1 September 2025.}\label{lst:g12:relational-databases-sql:overdue}

\omoutput[1-11,17-40]{code/grade-12/09-relational-databases-sql/outputs/overdue.txt}{The number of loans overdue, then
the loans themselves (the statements are in the listing).}\label{out:g12:relational-databases-sql:overdue}

\omoutput[49-68,74-80]{code/grade-12/09-relational-databases-sql/outputs/overdue.txt}{The members concerned, with their
streets, and the days of lateness.}\label{out:g12:relational-databases-sql:overdue-2}

\needspace{4\baselineskip}\medskip\noindent\textbf{Part I --- The schema.}
\begin{enumerate}
  \item Which \omterm{def:g12:relational-databases-sql:relation}{relations} does the problem need? Draw them with their keys and \omterm{def:g12:relational-databases-sql:keys}{foreign keys}.
  \item Why is the date of return allowed to be \code{NULL}, but not the date of loan?
  \item Why does \code{loans} store \code{member_id} and not the member's name?
  \item What would go wrong in the last \omterm{def:g12:relational-databases-sql:sql}{query} if a member could appear in \code{loans} without existing in
    \code{members}?
\end{enumerate}

\needspace{4\baselineskip}\medskip\noindent\textbf{Part II --- Overdue on a given day.}
\begin{enumerate}[resume]
  \item Explain each of the three parts of the condition: \code{out_date <= D}, \code{due_date < D}, and the part in
    brackets.
  \item Why are the brackets necessary?
  \item How many loans were overdue on 1 September? How many of them were returned later, and when?
  \item Why does \code{ORDER BY l.due_date} put the longest-overdue loans first?
\end{enumerate}

\needspace{4\baselineskip}\medskip\noindent\textbf{Part III --- Members and streets.}
\begin{enumerate}[resume]
  \item The third \omterm{def:g12:relational-databases-sql:sql}{query} joins four \omterm{def:g12:relational-databases-sql:relation}{relations}. Follow the path of the joins from a loan to a street.
  \item There are $18$ loans but $14$ members. Why fewer, and what does \code{DISTINCT} do here?
  \item Which streets appear twice, and is that a \omterm{def:g12:relational-databases-sql:anomaly}{redundancy} in the \omterm{def:g12:relational-databases-sql:database}{database}?
  \item Without \code{DISTINCT}, how many \omterm{def:g12:relational-databases-sql:relation}{rows} would the \omterm{def:g12:relational-databases-sql:sql}{query} return?
  \item A member moves to another street. How many \omterm{def:g12:relational-databases-sql:relation}{rows} of the \omterm{def:g12:relational-databases-sql:database}{database} change, and how many \omterm{def:g12:relational-databases-sql:relation}{rows} of this answer?
\end{enumerate}

\needspace{4\baselineskip}\medskip\noindent\textbf{Part IV --- Days of lateness.}
\begin{enumerate}[resume]
  \item What does \code{julianday(D) - julianday(due_date)} measure?
  \item How many days of lateness had piled up on 1 September, in all? The longest single delay?
  \item Which loan has the longest delay, and when was it due?
  \item The average delay is not in the output. Compute it from the output, then write the \omterm{def:g12:relational-databases-sql:sql}{query}.
  \item For the four loans returned in September, how many days late were they on the day they came back?
  \item The library fines $10$ cents per day late. What did the overdue loans owe on 1 September?
  \item State the result: the number of loans overdue on 1 September 2025, the members concerned, and the total days
    of lateness.
\end{enumerate}
\end{problem}
